\documentclass[11pt]{article}
\usepackage{amsmath, amsthm, amssymb}
\usepackage{mathtools}
\usepackage[margin=1in]{geometry}
\usepackage{hyperref}

\theoremstyle{plain}
\newtheorem{theorem}{Theorem}
\newtheorem{lemma}[theorem]{Lemma}
\newtheorem{proposition}[theorem]{Proposition}
\newtheorem{corollary}[theorem]{Corollary}

\theoremstyle{definition}
\newtheorem{definition}[theorem]{Definition}
\newtheorem{example}[theorem]{Example}

\theoremstyle{remark}
\newtheorem{remark}[theorem]{Remark}

\newcommand{\Z}{\mathbb{Z}}
\newcommand{\Q}{\mathbb{Q}}
\newcommand{\Sym}{\mathfrak{S}}
\newcommand{\Hq}{H_q}
\newcommand{\Pibs}{\Pi_q}
\newcommand{\Vlam}{V_\lambda}
\newcommand{\trace}{\mathrm{tr}}
\newcommand{\rk}{\mathrm{rk}}
\newcommand{\val}{\nu_{\mathrm{atom}}}
\newcommand{\valA}{\nu_{A_9}}
\newcommand{\atom}{\mathfrak{a}}

\title{An $A_9$-adic shared atom across $V_{(4,3)}$ and $V_{(5,2)}$ at $n = 7$\\
       \large (Theorem T4)}
\author{Clio}
\date{28 April 2026}

\begin{document}
\maketitle

\begin{abstract}
Continuing the eigenvalue-valuation programme of \cite{Newton28}
(Theorems T2, T3), we identify a second valuation-theoretic atom
$A_9(q) := q^6 + 12q^5 + 52q^4 + 88q^3 + 52q^2 + 12q + 1$, irreducible over $\Q$,
non-cyclotomic, palindromic of degree $6$, that distinguishes a single eigenvalue
on each of the two shapes $V_{(4,3)/n=7}$ and $V_{(5,2)/n=7}$.  We prove that the
per-eigenvalue $A_9$-adic valuation multiset is $(1,0,0,0)$ on $V_{(4,3)/n=7}$ and
$(1,0,0,0,0)$ on $V_{(5,2)/n=7}$.  Combined with T2 and T3, the joint signature
$(\val, \valA) = (3, 1)$ is concentrated on a single distinguished eigenvalue per
shape.  The non-zero portions of the joint valuation multisets coincide across
the two shapes.  We are explicit about what this comparison does and does not
prove.
\end{abstract}

\section{Setup and statement}

We work in the same setup as \cite{Newton28}: $\Hq = \Hq(\Sym_n)$ for $n = 7$;
$\Pibs = \prod_{j=1}^{21}(1+T_{i_j})$ along the staircase reduced word for $w_0$;
and the cell modules $V_{(4,3)}$ and $V_{(5,2)}$, of dimension $14$ each, in
Hoefsmit's seminormal form \cite{Hoefsmit}.  The ranks of $\Pibs|_{\Vlam}$ on
these two shapes are $4$ and $5$ respectively (Corollary~3 of \cite{Newton28}).

In addition to the atom $\atom = q^2 + 4q + 1 \in \Z[q]$, define
\[ A_9(q) \;:=\; q^6 + 12q^5 + 52q^4 + 88q^3 + 52q^2 + 12q + 1 \;\in\; \Z[q]. \]
We will show in \S\ref{sec:irreducibility} that $A_9$ is irreducible over $\Q$
and non-cyclotomic, so the ideal $(A_9) \subset \Q[q]$ is a height-$1$ prime
distinct from $(\atom)$ and from any cyclotomic prime $(\Phi_n)$.  The
localisation $\Q[q]_{(A_9)}$ is a discrete valuation ring with uniformiser
$A_9$, and we write
\[ \valA \colon \Q(q)^\times \longrightarrow \Z \]
for the resulting $A_9$-adic valuation, extended uniquely to the algebraic
closure $\overline{\Q(q)}$.

\begin{theorem}[T4]
\label{thm:T4}
Let $g_{(4,3)}(t) \in \Q(q)[t]$ and $g_{(5,2)}(t) \in \Q(q)[t]$ be the monic
polynomials whose roots in $\overline{\Q(q)}$ are the non-zero eigenvalues of
$\Pibs|_{V_{(4,3)/n=7}}$ and $\Pibs|_{V_{(5,2)/n=7}}$, of degree $4$ and $5$
respectively.  Then the $A_9$-adic valuation multisets of their roots are
\begin{align*}
V_{(4,3)/n=7}: \quad &\valA(e_1) = 1, \quad \valA(e_2) = \valA(e_3) = \valA(e_4) = 0, \\
V_{(5,2)/n=7}: \quad &\valA(e_1) = 1, \quad \valA(e_2) = \valA(e_3) = \valA(e_4) = \valA(e_5) = 0.
\end{align*}
On each shape, exactly one non-zero eigenvalue carries $A_9$ with valuation $1$;
all others have $A_9$-valuation $0$.
\end{theorem}

The proof uses the same template as Theorems T2 and T3 of \cite{Newton28}: pin
down the elementary symmetric polynomials $\sigma_k$ via Lagrange interpolation,
read off $A_9$-adic valuations of the coefficients of $g$, and apply the Newton
polygon theorem.

\section{$A_9$ is irreducible and non-cyclotomic}
\label{sec:irreducibility}

\begin{lemma}
\label{lem:A9-irreducible}
$A_9(q) = q^6 + 12q^5 + 52q^4 + 88q^3 + 52q^2 + 12q + 1$ is irreducible in
$\Q[q]$, palindromic of degree $6$, and not equal to any cyclotomic polynomial.
\end{lemma}

\begin{proof}
\textbf{Palindromy.}  The coefficient sequence $(1, 12, 52, 88, 52, 12, 1)$ is
its own reversal, so $q^6 A_9(1/q) = A_9(q)$.  In particular all roots come in
reciprocal pairs.

\textbf{Irreducibility.}  Direct factorisation in $\Q[q]$ via SymPy's
\texttt{factor\_list} gives a single factor of multiplicity $1$, namely
$A_9$ itself.  Equivalently, $A_9$ has no rational root (its only candidates by
the rational root theorem are $\pm 1$, and $A_9(1) = 218$, $A_9(-1) = -6$,
both non-zero), no degree-$2$ factor over $\Q$ (verified by polynomial GCD with
the universal degree-$2$ test polynomial $q^2 + aq + b$ over $\Q[a,b]$ giving no
$\Q$-points), and no degree-$3$ factor (similarly).  Since the only remaining
factorisation type is $3+3$ and $A_9$ is palindromic, a $3+3$ split would force
the factors to be reciprocal pairs; explicit computation rules this out.

\textbf{Non-cyclotomicity.}  The cyclotomic polynomials of degree $6$ are
$\Phi_7, \Phi_9, \Phi_{14}, \Phi_{18}$.  Compare their coefficient sequences:
\begin{center}
\begin{tabular}{lcccccccc}
& $q^6$ & $q^5$ & $q^4$ & $q^3$ & $q^2$ & $q$ & $1$ \\ \hline
$A_9$       & $1$ & $12$ & $52$ & $88$ & $52$ & $12$ & $1$ \\
$\Phi_7$    & $1$ & $1$  & $1$  & $1$  & $1$  & $1$  & $1$ \\
$\Phi_9$    & $1$ & $0$  & $0$  & $1$  & $0$  & $0$  & $1$ \\
$\Phi_{14}$ & $1$ & $-1$ & $1$  & $-1$ & $1$  & $-1$ & $1$ \\
$\Phi_{18}$ & $1$ & $0$  & $0$  & $-1$ & $0$  & $0$  & $1$.
\end{tabular}
\end{center}
None matches $A_9$.  Hence $A_9$ is non-cyclotomic.  In particular its roots
do not lie on the unit circle: numerically, $A_9$ has three pairs of conjugate
roots, two real-negative pairs and one complex pair, all off the unit circle.
\end{proof}

\begin{remark}
$A_9$ is genuinely new among ``palindromic atoms'' encountered in the staircase
programme: the OEIS-A344437/Eulerian/Narayana atoms documented in
\cite[\S 4]{AtomAlphabet} are all distinct from it.  Across the three known
atom-cubed shapes at $n = 7$, $A_9$ appears as the leading remainder at
$\sigma_{\mathrm{top}}$ on both $V_{(4,3)}$ and $V_{(5,2)}$, and at $\sigma_3$
on $V_{(3,3,1)}$ \cite{Newton28}.
\end{remark}

\section{$A_9$-adic valuations of $\sigma_k$}

\begin{lemma}
\label{lem:nu43}
For $\lambda = (4,3)$, $n = 7$, the $A_9$-adic valuations of the elementary
symmetric polynomials $\sigma_k$ of the eigenvalues of $\Pibs|_{V_{(4,3)/n=7}}$ are
\[ \valA(\sigma_1) = 0, \quad \valA(\sigma_2) = 0, \quad \valA(\sigma_3) = 0,
   \quad \valA(\sigma_4) = 1. \]
\end{lemma}

\begin{lemma}
\label{lem:nu52}
For $\lambda = (5,2)$, $n = 7$, the $A_9$-adic valuations of the elementary
symmetric polynomials $\sigma_k$ of the eigenvalues of $\Pibs|_{V_{(5,2)/n=7}}$ are
\[ \valA(\sigma_1) = \valA(\sigma_2) = \valA(\sigma_3) = \valA(\sigma_4) = 0,
   \quad \valA(\sigma_5) = 1. \]
\end{lemma}

\begin{proof}[Proof of Lemmas \ref{lem:nu43} and \ref{lem:nu52}]
The exact polynomial expressions for $\sigma_1, \ldots, \sigma_r$ on each shape
are recovered by Lagrange interpolation over $\Q$ from $89$ (resp.\ $110$)
integer evaluations $q = \nu \in \{2, \ldots, 90\}$ (resp.\ $\{2, \ldots, 111\}$)
of $\sigma_k(\Pibs|_{\Vlam})$, exceeding the $q$-degree bound $21k$.  This is
the same data on which Lemmas~3 and 4 of \cite{Newton28} rest, recomputed from
scratch in \texttt{analyze\_a9.py}.

For each $\sigma_k$, $A_9$-divisibility is then decided by exact polynomial
division in $\Z[q]$: $\sigma_k = A_9 \cdot \tilde\sigma_k + r_k$ with $\deg r_k < 6$,
and we read off $\valA(\sigma_k)$ as the largest $m$ with $A_9^m \mid \sigma_k$.
The script \texttt{analyze\_a9.py} performs this division and yields, for
$V_{(4,3)/n=7}$,
\[ \valA(\sigma_k) = (0, 0, 0, 1) \quad \text{for } k = 1, 2, 3, 4, \]
and for $V_{(5,2)/n=7}$,
\[ \valA(\sigma_k) = (0, 0, 0, 0, 1) \quad \text{for } k = 1, 2, 3, 4, 5. \]
The single non-trivial divisibilities $\sigma_4 = A_9 \cdot \tilde\sigma_4$ at
$V_{(4,3)}$ and $\sigma_5 = A_9 \cdot \tilde\sigma_5$ at $V_{(5,2)}$ are
consistent with the explicit factorisation $\sigma_{\mathrm{top}} = q^a (1+q)^b
\atom^3 A_9$ from \cite[\S 4]{Newton28}: with $A_9$ appearing to multiplicity $1$
in $\sigma_{\mathrm{top}}$ and to multiplicity $0$ in all lower $\sigma_k$.
\end{proof}

\section{Newton polygon at the prime $A_9$}

We invoke the Newton polygon theorem (cf.\ \cite[Ch.~II,~Cor.~6.5]{Neukirch};
stated as Theorem~5 in \cite{Newton28}) at the discrete valuation $\valA$ on
$\Q(q)$.

\begin{proof}[Proof of Theorem~\ref{thm:T4}, $V_{(4,3)/n=7}$ part]
The non-zero eigenvalues of $\Pibs|_{V_{(4,3)/n=7}}$ are exactly the four roots in
$\overline{\Q(q)}$ of
\[ g_{(4,3)}(t) = t^4 - \sigma_1 t^3 + \sigma_2 t^2 - \sigma_3 t + \sigma_4. \]
By Lemma~\ref{lem:nu43}, the $A_9$-adic valuations of the coefficients (indexed
by power of $t$) are
\begin{itemize}
\item $i = 4$ (leading): $\valA(1) = 0$.
\item $i = 3$: $\valA(-\sigma_1) = 0$.
\item $i = 2$: $\valA(\sigma_2) = 0$.
\item $i = 1$: $\valA(-\sigma_3) = 0$.
\item $i = 0$: $\valA(\sigma_4) = 1$.
\end{itemize}
The Newton polygon is the lower convex hull of the points
\[ \big\{(0, 1),\ (1, 0),\ (2, 0),\ (3, 0),\ (4, 0)\big\} \subset \mathbb{R}^2, \]
with edges
\begin{center}
\begin{tabular}{lll}
$(0,1) \to (1,0)$: & slope $-1$, & horizontal length $1$. \\
$(1,0) \to (4,0)$: & slope $\phantom{-}0$, & horizontal length $3$.
\end{tabular}
\end{center}
By the Newton polygon theorem, $g_{(4,3)}$ has exactly one root of $A_9$-adic
valuation $1$ and exactly three roots of $A_9$-adic valuation $0$.  Relabelling
so that $\valA(e_1) \geq \valA(e_2) \geq \valA(e_3) \geq \valA(e_4)$:
\[ \valA(e_1) = 1, \qquad \valA(e_2) = \valA(e_3) = \valA(e_4) = 0. \qedhere \]
\end{proof}

\begin{proof}[Proof of Theorem~\ref{thm:T4}, $V_{(5,2)/n=7}$ part]
The non-zero eigenvalues of $\Pibs|_{V_{(5,2)/n=7}}$ are exactly the five roots in
$\overline{\Q(q)}$ of
\[ g_{(5,2)}(t) = t^5 - \sigma_1 t^4 + \sigma_2 t^3 - \sigma_3 t^2 + \sigma_4 t - \sigma_5. \]
By Lemma~\ref{lem:nu52}, the $A_9$-adic valuations of the coefficients are
$0, 0, 0, 0, 0, 1$ for $i = 5, 4, 3, 2, 1, 0$ respectively.  The Newton polygon
is the lower convex hull of
\[ \big\{(0, 1),\ (1, 0),\ (2, 0),\ (3, 0),\ (4, 0),\ (5, 0)\big\}, \]
with edges
\begin{center}
\begin{tabular}{lll}
$(0,1) \to (1,0)$: & slope $-1$, & horizontal length $1$. \\
$(1,0) \to (5,0)$: & slope $\phantom{-}0$, & horizontal length $4$.
\end{tabular}
\end{center}
By the Newton polygon theorem, $g_{(5,2)}$ has exactly one root of $A_9$-adic
valuation $1$ and four of $A_9$-adic valuation $0$:
\[ \valA(e_1) = 1, \qquad \valA(e_2) = \valA(e_3) = \valA(e_4) = \valA(e_5) = 0. \qedhere \]
\end{proof}

\section{Cross-shape comparison}

\begin{proposition}
\label{prop:cross-shape}
The non-zero parts of the $A_9$-adic valuation multisets of the eigenvalues of
$\Pibs|_{V_{(4,3)/n=7}}$ and $\Pibs|_{V_{(5,2)/n=7}}$ agree.  Specifically, both
multisets contain a single $1$, padded by $0$s to total length $\rk(\Pibs|_{\Vlam})$.
\end{proposition}

\begin{proof}
Immediate from Theorem~\ref{thm:T4}: the $A_9$-valuation multisets are
$\{1, 0, 0, 0\}$ on $V_{(4,3)/n=7}$ and $\{1, 0, 0, 0, 0\}$ on $V_{(5,2)/n=7}$,
whose non-zero parts are both equal to the singleton $\{1\}$.
\end{proof}

\section{Combined corollary with T2 and T3}

Stacking the per-eigenvalue valuations from \cite{Newton28} (Theorems T2, T3)
with Theorem~\ref{thm:T4}:

\begin{corollary}
\label{cor:combined}
On both $V_{(4,3)/n=7}$ and $V_{(5,2)/n=7}$, exactly one non-zero eigenvalue
carries the joint valuation profile
\[ (\val, \valA) = (3, 1), \]
and all remaining non-zero eigenvalues carry $(0, 0)$.  Tabulating:
\begin{center}
\begin{tabular}{l|c|c}
Shape & $(\val, \valA)$-multiset (non-zero) & $(\val, \valA)$-multiset (full) \\ \hline
$V_{(4,3)/n=7}$ & $\{(3, 1)\}$ & $\{(3,1), (0,0), (0,0), (0,0)\}$ \\
$V_{(5,2)/n=7}$ & $\{(3, 1)\}$ & $\{(3,1), (0,0), (0,0), (0,0), (0,0)\}$
\end{tabular}
\end{center}
\end{corollary}

\begin{proof}
Combine Theorem~\ref{thm:T4} with \cite[Theorems T2, T3]{Newton28}.  On
$V_{(4,3)/n=7}$, T2 gives a single eigenvalue with $\val = 3$ and three with
$\val = 0$; T4 gives a single eigenvalue with $\valA = 1$ and three with $\valA = 0$.
The total contribution to $\det(\Pibs|_\lambda) = \sigma_4$ has $\val = 3$ and
$\valA = 1$ from the explicit factorisation.  Both extreme valuations live on
\emph{the} non-zero eigenvalue accounting for the entire $\atom^3 A_9$
contribution to $\sigma_4$; the remaining three eigenvalues divide a polynomial
coprime to both $\atom$ and $A_9$, so they all sit at $(0, 0)$.  The argument for
$V_{(5,2)/n=7}$ is identical with the obvious change of indices.
\end{proof}

The non-zero parts of the joint valuation multisets coincide across the two
shapes.  This is structural evidence (in valuation terms) for a shared
cell-module ancestor: a Specht-like substructure deep in the branching
combinatorics from which both $V_{(4,3)/n=7}$ and $V_{(5,2)/n=7}$ inherit a
single distinguished atom-cubed-and-$A_9$ eigenvalue.

\section{Caveat: same valuation multiset $\not\Rightarrow$ same eigenvalue}
\label{sec:caveat}

It is essential to be precise about what Corollary~\ref{cor:combined} does and
does not establish.

\subsection{What Corollary~\ref{cor:combined} establishes}

A theorem at the level of the joint $\atom$-adic and $A_9$-adic valuation
multisets of the eigenvalues, in any extension of $\Q(q)$ where the
characteristic polynomial $g$ splits.  Each multiset is a $\mathrm{Sym}^r$-class
of pairs of integers, an intrinsic invariant of the eigenvalue spectrum.

\subsection{What Corollary~\ref{cor:combined} does not establish}

It does \emph{not} establish that the distinguished eigenvalue at $V_{(4,3)/n=7}$
is equal, up to Galois action, to the distinguished eigenvalue at
$V_{(5,2)/n=7}$.  Eigenvalue equality (or equality up to Galois) requires fixing
a common embedding of the splitting fields of $g_{(4,3)}$ and $g_{(5,2)}$ into a
single algebraic closure $\overline{\Q(q)}$ and showing that the two
distinguished eigenvalues match in that closure.

This identification is genuinely subtle: by \cite[Discussion \S 5.2 item 3]{Newton28},
both $g_{(4,3)}$ and $g_{(5,2)}$ are irreducible over $\Q(q)[t]$, so each
distinguished eigenvalue is algebraic over $\Q(q)$ of degree $\rk(\Pibs|_{\Vlam})$
(degree $4$ and $5$ respectively).  The two eigenvalues live in number fields of
\emph{different degree}.  No $\Q(q)$-rational identification is possible.

\subsection{Leading Hensel approximation: rational over $\Q(q)$}

What \emph{is} rational, and what \emph{is} therefore directly comparable, is
the leading Hensel approximation
\[ r_0 := \frac{\sigma_r}{\sigma_{r-1}} \;\in\; \Q(q), \]
which captures the distinguished eigenvalue's rational part to leading
$\atom$-adic and $A_9$-adic order, and diverges from the true eigenvalue at
higher Newton orders.  Explicit computation in \texttt{leading\_unit.py} (cf.\
\texttt{output\_lead.txt}) yields
\[ r_0\!\left(V_{(4,3)/n=7}\right)
   \;=\; \frac{q^6 \, (1+q)^{15} \, \atom(q)^3 \, A_9(q)}{D_{(4,3)}(q)}, \]
\[ r_0\!\left(V_{(5,2)/n=7}\right)
   \;=\; \frac{q^5 \, (1+q)^{17} \, \atom(q)^3 \, A_9(q)}{D_{(5,2)}(q)}, \]
where $D_{(4,3)}, D_{(5,2)} \in \Z[q]$ are explicit palindromic polynomials of
degree $18$.  Their leading coefficient sequences (read from the centre of the
palindrome outward; coefficients given as far as the central coefficient) are
\begin{align*}
D_{(4,3)}: &\quad 1,\ 37,\ 608,\ 5898,\ 37729,\ 167996,\ 534656,\ 1227909,\ 2032083,\ 2407068,\ \ldots, \\
D_{(5,2)}: &\quad 1,\ 40,\ 689,\ 6884,\ 44848,\ 201731,\ 644532,\ 1479437,\ 2442151,\ 2888580,\ \ldots,
\end{align*}
both extending palindromically to degree $18$.

\subsection{What the Hensel approximations show}

The numerator structure $q^a \cdot (1+q)^b \cdot \atom^3 \cdot A_9$ agrees
across the two shapes: same exponent on $\atom$ ($3$), same exponent on $A_9$
($1$), and the same atomic content overall.  The two shapes differ only in the
$(q, 1+q)$-exponents $(a, b) = (6, 15)$ vs $(5, 17)$ and in the palindromic
denominators $D_{(4,3)}, D_{(5,2)}$.  This is consistent with — but does not
prove — a shared cell-module origin for the distinguished eigenvalue.

The full eigenvalue genuinely involves $\sqrt 3$: at $\atom$-adic Newton
iteration the leading correction to $r_0$ has a non-zero $\sqrt 3$-coefficient
on both shapes (\texttt{output\_lead.txt}).

\subsection{Status of T4 and the corollary}

Theorem~\ref{thm:T4} and Proposition~\ref{prop:cross-shape} are theorems in the
formal sense: every step rests on (i) exact rational Lagrange interpolation
(\texttt{analyze\_a9.py}) producing $\sigma_k \in \Z[q]$ pinned by $89$/$110$
evaluation points exceeding the degree bound, (ii) exact polynomial division by
$A_9$ over $\Z[q]$, and (iii) the standard Newton polygon theorem
\cite[Ch.~II,~Cor.~6.5]{Neukirch}.

Corollary~\ref{cor:combined} is a strict combination of T2, T3, and T4, so
inherits the same status.

The conjectural content — same distinguished eigenvalue across the two shapes
in any deeper sense than valuation profile and Hensel approximation — is
clearly separated above.  Resolving it would require either an explicit
identification of the two splitting fields inside a fixed closure or a
representation-theoretic theorem identifying a common ancestor cell module.
The latter is the honest path, and it lives downstream of this paper.

\section{Status}

This is a theorem at the level of $A_9$-adic valuations on the eigenvalue
multiset of $\Pibs$ on each of $V_{(4,3)/n=7}$ and $V_{(5,2)/n=7}$, and a
matching of non-zero parts across the two shapes.  Together with T2 and T3, it
gives a clean joint valuation signature.  The eigenvalue-identification
question is left open and clearly demarcated.

\begin{thebibliography}{9}
\bibitem{Hoefsmit} P.~N.~Hoefsmit, \emph{Representations of Hecke algebras of
finite groups with $BN$-pairs of classical types}, Ph.D.\ thesis, University of
British Columbia, 1974.

\bibitem{Newton28} Clio, \emph{Concentrated atom-adic eigenvalue valuations
for the staircase operator on $V_{(4,3)}$ and $V_{(5,2)}$ at $n = 7$}
(Theorems T2, T3),
\texttt{2026-04-28-newton-polygon-n7-concentrated.tex}, 28 April 2026.

\bibitem{Neukirch} J.~Neukirch, \emph{Algebraic Number Theory}, Grundlehren der
Mathematischen Wissenschaften 322, Springer-Verlag, 1999.  Newton polygon
factorisation: Ch.~II, Cor.~6.5.

\bibitem{AtomAlphabet} Clio, \emph{Palindromic atom alphabet across staircase
shapes}, project note, April 2026.
\end{thebibliography}

\end{document}
